\subsection{Hodge numbers}
We only make some brief remarks on the Hodge numbers, and do not include any detailed calculations. 
The coherent Leray calculation is similar to
\cite[Proposition~9.13]{AC}. For this construction it gives
\begin{align}\label{direct}
 f_*\cO_X=\cO_{\PP^1},\qquad
 R^1f_*\cO_X\simeq\cO_{\PP^1}\oplus\cO_{\PP^1}(-1),\qquad
 R^2f_*\cO_X\simeq\cO_{\PP^1}(-1).
\end{align}
Here we use the cohomology of the non-reduced fiber $4B_*$, rather
than that of its reduction $B_*$.
Each fiber $D$ has
$h^q(D,\cO_D)=(1,2,1)$ for $q=0,1,2$, so Grauert's theorem gives
local freeness and base change. In other words, $f$ is cohomologically
flat in all degrees. The map $f$ is proper and flat, and its fibers
are Cartier divisors in the smooth threefold $X$, hence Gorenstein.
Analytic relative duality of Ramis--Ruget--Verdier \cite{RRV} therefore
gives
\begin{align*}
 (R^2f_*\cO_X)^*\simeq f_*\omega_{X/\PP^1},\qquad
 \omega_{X/\PP^1}=K_X\otimes f^*K_{\PP^1}^{-1}
 \simeq f^*\cO_{\PP^1}(1).
\end{align*}
Since $f_*\cO_X=\cO_{\PP^1}$, the projection formula yields
$R^2f_*\cO_X\simeq\cO_{\PP^1}(-1)$.

The invariant real class $\psi$ gives a
nowhere-zero section of $R^1f_*\cO_X$: this follows from the
normalization sequences at the three reduced singular fibers and by
restriction to $B_*$ at the multiple fiber. Riemann--Roch gives
$\deg R^1f_*\cO_X=-1$, and the resulting extension of $\cO(-1)$ by
$\cO$ splits because $H^1(\PP^1,\cO(1))=0$.
The coherent Leray spectral sequence then yields
\begin{align}\label{hodge}
 h^{0,0}(X)=1,\qquad h^{0,1}(X)=1,\qquad
 h^{0,2}(X)=h^{0,3}(X)=0.
\end{align}
We do not compute the remaining Hodge numbers.

